\chapter{Expressions algébriques, équations et inéquations}
\label{workbook:chapter:02}
\section{Variables, monômes et polynômes}
\label{workbook:section:02.01}
\begin{exerciseblock}{Variables, monômes et polynômes}
\item Pour $P(x,y)=3x^2y-5xy^2+7$, préciser le degré de chaque terme, le degré total de $P$, puis calculer $P(2,-1)$.
\item Réduire et ordonner : $4x^2-3x+7-(2x^2+5x-1)+3(x^2-2)$.
\item Un rectangle a pour côtés $x+3$ et $2x-1$. Écrire son aire et son périmètre sous forme développée, puis préciser les valeurs de $x$ admissibles dans le contexte.
\end{exerciseblock}
\section{Factorisation et fractions algébriques}
\label{workbook:section:02.02}
\begin{exerciseblock}{Factorisation et fractions algébriques}
\item Factoriser complètement : $6x^3-24x$, $x^2-10x+25$, $4x^2-9$.
\item Simplifier en indiquant les restrictions : $\dfrac{x^2-9}{x^2+x-6}$.
\item Comparer les fonctions $f(x)=\dfrac{x^2-1}{x-1}$ et $g(x)=x+1$. Sont-elles égales comme fonctions sur $\R$?
\end{exerciseblock}
\section{Équations du premier degré}
\label{workbook:section:02.03}
\begin{exerciseblock}{Équations du premier degré}
\item Résoudre $4(2x-3)-5=3(x+7)$ et vérifier la solution par substitution.
\item Résoudre $\dfrac{x-2}{3}-\dfrac{2x+1}{4}=5$ en éliminant les dénominateurs.
\item Deux forfaits coûtent $A(n)=18+0{,}08n$ et $B(n)=27+0{,}03n$. Pour quel nombre de messages les coûts sont-ils égaux? Interpréter.
\end{exerciseblock}
\section{Inéquations et tableaux de signes}
\label{workbook:section:02.04}
\begin{exerciseblock}{Inéquations et tableaux de signes}
\item Résoudre $-3(2x-1)>5x+7$ et expliquer le changement de sens lors de la dernière étape.
\item Construire un tableau de signes et résoudre $(x-4)(2x+1)\le0$.
\item Résoudre $\dfrac{x+2}{x-3}>0$ en distinguant zéros et valeurs interdites.
\end{exerciseblock}
\section{Équations avec racines ou valeurs absolues}
\label{workbook:section:02.05}
\begin{exerciseblock}{Équations avec racines ou valeurs absolues}
\item Résoudre $\sqrt{2x+3}=x$ en indiquant d'abord les contraintes de domaine.
\item Résoudre $|3x-5|=7$ puis interpréter les deux solutions comme deux distances égales.
\item Résoudre $|x+1|<4$ et $|2x-3|\ge5$ sous forme d'intervalles.
\end{exerciseblock}
\section{Équivalences et transformations à contrôler}
\label{workbook:section:02.06}
\begin{exerciseblock}{Équivalences et transformations à contrôler}
\item Pour chacune des étapes suivantes, préciser s'il s'agit d'une équivalence : ajouter $5$ aux deux membres; multiplier par $x$; élever au carré; prendre la racine carrée.
\item Résoudre $x=\sqrt{x+6}$ en conservant les solutions candidates puis en les vérifiant.
\item Expliquer pourquoi l'étape $x^2=9\Rightarrow x=3$ perd une solution et écrire l'équivalence correcte.
\end{exerciseblock}
\section{Équivalence, implication et ensemble des solutions}
\label{workbook:section:02.07}
\begin{exerciseblock}{Équivalence, implication et ensemble des solutions}
\item Déterminer les ensembles de solutions de $x^2=4$ et $x=2$, puis décider quelle implication est vraie.
\item Donner une condition suffisante mais non nécessaire pour qu'un entier soit divisible par $6$, puis une condition nécessaire mais non suffisante.
\item Réécrire la résolution de $2x+5=11$ comme une chaîne d'équivalences complète.
\end{exerciseblock}
\section{Résoudre sans perdre d'information}
\label{workbook:section:02.08}
\begin{exerciseblock}{Résoudre sans perdre d'information}
\item Résoudre $\dfrac{x+1}{x-2}=3$ en écrivant d'abord le domaine et en vérifiant la solution.
\item Résoudre $\sqrt{x+1}=x-1$ en indiquant toutes les contraintes avant de mettre au carré.
\item Concevoir une procédure de contrôle en quatre étapes pour une équation comportant à la fois un dénominateur et une racine.
\end{exerciseblock}
\section*{Synthèse du chapitre}
\begin{synthesisbox}
\item Résoudre et comparer les méthodes : $x^2-5x+6=0$ par factorisation puis par formule quadratique.
\item Une entreprise facture $C(x)=120+18x$ et reçoit $R(x)=30x$. Déterminer le seuil de rentabilité et interpréter l'inéquation $R(x)>C(x)$.
\item Résoudre $\dfrac{x-1}{x+2}\le2$ à l'aide d'un tableau de signes, puis contrôler avec deux valeurs tests.
\end{synthesisbox}
